\section{The endomorphism algebra of the residual complex}\label{sec:I}
The finite topology on $\Hom_A(M,N)$ means pointwise convergence
when the modules are discrete. The universal problem below uses
separated complete graded topological $A$-algebras, continuous
multiplication, and continuous graded homomorphisms, with $A$ discrete
as base. The prescribed degree zero has the product of local adic
topologies. This specifies an interpretation of completion; the
uncertain transcription on 43:65 does not uniquely specify a category.

\begin{proposition}[I1: degree-zero diagonality; Proved]\label{i:zero}
For a regular noetherian local ring, normalize the residual complex by
$K^i=\bigoplus_{\operatorname{ht}\mathfrak p=i}E_A(A/\mathfrak p)$.
Then $\End_A^0(K)=\prod_{\mathfrak p}\widehat{A_{\mathfrak p}}$
as topological algebras, and $\End_A^i(K)=0$ for $i<0$.
\end{proposition}
\begin{proof}
Write $E_{\mathfrak p}=E_A(A/\mathfrak p)$. If
$a\in\mathfrak p\setminus\mathfrak q$, every element of
$E_{\mathfrak p}$ is killed by a power of $a$, while multiplication
by $a$ is invertible on $E_{\mathfrak q}$. Thus the corresponding Hom
vanishes unless $\mathfrak p\subseteq\mathfrak q$. Distinct primes
of the same height are incomparable; projections detect every element
of a direct sum. This proves diagonality and negative-degree vanishing.
Matlis' endomorphism theorem gives the displayed rings
\cite[Lemma 2, proof, p.~123]{Sastry1992}. Annihilators of finite
subsets of $E_{\mathfrak p}$ give exactly the local adic topology:
the submodules killed by powers of the maximal ideal have finite
length over $A_{\mathfrak p}$ and exhaust the hull. The direct sum
gives the product topology.
\end{proof}

\begin{theorem}[I2: complete surface presentation; Proved]\label{i:surface}
Let $A=k[[s,t]]$, $F=\operatorname{Frac}(A)$, and
$\mathfrak m=(s,t)$. Then $\End_A^*(K)$ with differential $\theta$
is universal among the separated complete graded topological
$A$-algebras with degree zero
\[
 B^0=F\times\prod_{\operatorname{ht}\mathfrak p=1}
                 \widehat{A_{\mathfrak p}}\times A,
\]
one specified element of degree one, its square zero, and the
immediate-specialization incidence relations. In particular the
canonical map from the algebra defined by this precise universal
property is an isomorphism, including degree two.
\end{theorem}
\begin{proof}
Let $\mathcal P_1$ denote the height-one primes, put
$R_{\mathfrak p}=\widehat{A_{\mathfrak p}}$, choose a generator
$f_{\mathfrak p}$ of $\mathfrak p$, and put
$F_{\mathfrak p}=\operatorname{Frac}(R_{\mathfrak p})$.
The factorial Gorenstein ring $A$ has residual resolution
\begin{equation}\label{i:res}
 0\to A\to F\xrightarrow{d^0}
 I=\bigoplus_{\mathfrak p\in\mathcal P_1}E_{\mathfrak p}
 \xrightarrow{d^1}E\to0,
\end{equation}
where $E=E_A(k)$, $E_{\mathfrak p}=F/A_{\mathfrak p}$, and
$d^0$ consists of quotient maps. Write $\delta_{\mathfrak p}$ for
the components of $d^1$ and $D(-)=\Hom_A(-,E)$.

\emph{Individual arrows.}
By \cite[Theorem 1.1, p.~1]{Schenzel2013}, $D(E_{\mathfrak p})$
is free of rank one over $R_{\mathfrak p}$: $A/\mathfrak p$ is
complete of dimension one. Its element $\delta_{\mathfrak p}$ is
a basis. Indeed, applying $\Hom_A(A/\mathfrak p,-)$ to
\eqref{i:res} leaves only
$\kappa(\mathfrak p)\to\Hom_A(A/\mathfrak p,E)$.
The free resolution of $A/(f_{\mathfrak p})$ shows this map is
surjective with kernel isomorphic to $A/\mathfrak p$, so it is
nonzero. Dualizing
$0\to\kappa(\mathfrak p)\to E_{\mathfrak p}
\xrightarrow{f_{\mathfrak p}}E_{\mathfrak p}\to0$
shows that the kernel of restriction of $D(E_{\mathfrak p})$ to
$D(\kappa(\mathfrak p))$ is
$f_{\mathfrak p}D(E_{\mathfrak p})$. Hence
$\delta_{\mathfrak p}$ is a unit times a basis. We use it to identify
\begin{equation}\label{i:arrows}
 \Hom_A(E_{\mathfrak p},E)=R_{\mathfrak p},\qquad
 \Hom_A(F,E_{\mathfrak p})=F_{\mathfrak p}.
\end{equation}
For the second identity, apply $\Hom_{A_{\mathfrak p}}(-,E_{\mathfrak p})$
to $0\to A_{\mathfrak p}\to F\to E_{\mathfrak p}\to0$
and localize at $f_{\mathfrak p}$. The middle Hom is already an
$F$-vector space. Explicitly $a\in F_{\mathfrak p}$ sends
$f\in F$ to $fa\bmod R_{\mathfrak p}$, using
$E_{\mathfrak p}=F_{\mathfrak p}/R_{\mathfrak p}$.

\emph{All graded pieces.}
Put
\[
 P=\prod_{\mathfrak p}R_{\mathfrak p},\qquad
 \mathbb A=\left\{(a_{\mathfrak p})\in\prod_{\mathfrak p} F_{\mathfrak p}:
 a_{\mathfrak p}\in R_{\mathfrak p}\text{ for almost all }\mathfrak p\right\}.
\]
For $S=A\setminus\{0\}$, factoriality gives $\mathbb A=S^{-1}P$:
clear its finitely many negative valuations by a product of powers of
$f_{\mathfrak p}$. Conversely every denominator has finitely many
prime divisors. A family $F\to E_{\mathfrak p}$ has values in the
direct sum precisely when its coefficients are integral almost
everywhere: evaluate at $1$ for necessity and use the finitely many
poles of each $f\in F$ for sufficiency. Consequently
\begin{equation}\label{i:pieces}
 \End_A^1(K)=\mathbb A\oplus P,\qquad
 \End_A^2(K)=\mathbb A/F.
\end{equation}
For the second equality dualize \eqref{i:res}, using $D(E)=A$:
\begin{equation}\label{i:dual}
 0\to A\to P\to D(F)\to E\to0.
\end{equation}
The first map is diagonal, since its components are
$a\delta_{\mathfrak p}$. Localizing kills $E$ and leaves $D(F)$
unchanged, proving the equality. For $c\in P$ representing $I\to E$
and $a\in\mathbb A$ representing $F\to I$, composition is
\begin{equation}\label{i:product}
 c\circ a=[(c_{\mathfrak p}a_{\mathfrak p})_{\mathfrak p}]
                         \in\mathbb A/F.
\end{equation}
All other products of positive-degree arrows vanish by endpoints;
degrees above two vanish. The differential is $(1,1)$, with square
$[1]=0$. Thus the degree-two relation is precisely diagonal $F$,
not an independent zero relation for each intermediate prime.

\emph{The topologies in degree one.}
The first summand has the usual restricted-product topology: $P$ with
its product adic topology is open. Indeed, vanishing at $1/u$ means
membership in $uP$, and these form a basis. The second summand has
the product of topologies $\tau_{\mathfrak p}$ with bases
\begin{equation}\label{i:mixed}
 uA+f_{\mathfrak p}^n R_{\mathfrak p}\ (u\notin\mathfrak p),
 \quad\text{or equivalently}\quad
 \mathfrak m^N A+f_{\mathfrak p}^n R_{\mathfrak p}.
\end{equation}
To prove this, exactness and factoriality give
\[
 \ker\delta_{\mathfrak p}
  =(A[1/f_{\mathfrak p}]+A_{\mathfrak p})/A_{\mathfrak p}.
\]
A fraction whose other components vanish has no other prime factors
in its denominator. Any finite set in $E_{\mathfrak p}$ lies in the
cyclic $A$-module generated by $1/(u f_{\mathfrak p}^n)$ for some
$u\notin\mathfrak p$. Its annihilator in the Hom space is exactly
$uA+f_{\mathfrak p}^nR_{\mathfrak p}$: modulo
$f_{\mathfrak p}^n$, the element $c/u$ must come from $A$.
Since $(u,f_{\mathfrak p}^n)$ is $\mathfrak m$-primary, some
$\mathfrak m^N$ is contained in it; conversely choose
$u\in\mathfrak m^N\setminus\mathfrak p$. This proves the
equivalent basis. At $\mathfrak p=(s)$, it is coefficientwise
$t$-adic convergence on $k((t))[[s]]$. Hence
$t^N\delta_{\mathfrak p}\to0$, although $t^N$ is always an
$s$-adic unit. The two occurrences of $P$ require different topologies.

\emph{Degree two.}
Write $P_\tau$ for the product with \eqref{i:mixed}. The subgroup
$H=\{h\in D(F):h(1)=0\}$ is open and topologically $P_\tau/A$.
For the quotient assertion, if $M'\subset M$ and $E$ is injective,
restriction $D(M)\to D(M')$ is an open surjection for the finite
topologies over a noetherian ring. The image of the annihilator of a
finitely generated $N\subset M$ is exactly the annihilator of
$M'\cap N$: extend a map from $M'$ by zero on $N$, and then use
injectivity. The intersection is finitely generated, proving openness.
Apply this to $F/A\subset I$; the kernel is $D(E)=A$.
Furthermore $D(F)=\bigcup_{u\ne0}u^{-1}H$, since every element of
$E$ is killed by a nonzero element of $A$, and multiplication by $u$
is a topological automorphism. This specifies the topology on
$\mathbb A/F$ without an unjustified ordinary-adelic quotient.

\emph{The universal morphism.}
Let $C$ have the prescribed degree zero and $T\in C^1$ satisfying
the relations. Write $e_\eta,e_P,e_m$ for the three block
idempotents. Finite sums of the individual height-one idempotents
converge to $e_P$, so continuity and incidence give
\begin{equation}\label{i:T}
 T=e_PTe_\eta+e_mTe_P.
\end{equation}
For $a=b/u\in\mathbb A$, $b\in P$, $u\ne0$, and $c\in P$, set
\begin{align}\label{i:normal}
 \alpha_C(a)&=bTe_\eta u_\eta^{-1},&
 \beta_C(c)&=e_mTc,\\
 \gamma_C([a])&=e_mTbTe_\eta u_\eta^{-1}.\nonumber
\end{align}
The elements of $P$ occupy the height-one block, and $u_\eta^{-1}$
means $1/u$ in the generic block. Clearing denominators and using
$A$-linearity proves independence of representations. For diagonal
$a=v/u\in F$, the last expression is
$v e_mTe_PTe_\eta u_\eta^{-1}=0$, by \eqref{i:T} and $T^2=0$.
Thus $\gamma_C$ descends to $\mathbb A/F$.
These formulas give an algebra homomorphism, the identity in degree
zero: the only composable positive-degree product is
$\beta_C(c)\alpha_C(a)=\gamma_C([ca])$, agreeing with
\eqref{i:product}. All others vanish by endpoints. The diagonal
actions agree directly, using $A$-linearity for the closed-point
action and the right $F$-action at the generic point.

We check continuity carefully. The map $\alpha_C$ is continuous
on the open subgroup $P$ with product adic topology, and hence on all
$u^{-1}P$. For $\beta_C$, given a zero-neighborhood $U$, choose
$V$ with $V+V\subset U$. Continuity from product adic degree zero
gives a finite set $T_0$ of primes and powers
$f_{\mathfrak p}^{n_{\mathfrak p}}$ such that tuples satisfying
those divisibilities map into $V$; all remaining coordinates are
unrestricted. Each map
$A\to C^1$, $a\mapsto a e_mTe_{\mathfrak p}$, is continuous
for the $\mathfrak m$-adic topology. Choose $N$ so that the sum of
its images of $\mathfrak m^N$ over $\mathfrak p\in T_0$ lies
in $V$. Tuples satisfying
$c_{\mathfrak p}\in\mathfrak m^N A+
f_{\mathfrak p}^{n_{\mathfrak p}}R_{\mathfrak p}$ split into a
tuple controlled by the first condition and finitely many such
$\mathfrak m^N A$ terms. Their images lie in $U$. This proves
continuity from $P_\tau$.
The continuous map $c\mapsto\beta_C(c)Te_\eta$ vanishes on
diagonal $A$, hence descends to the open quotient $H=P_\tau/A$.
Multiplication by $u_\eta^{-1}$ gives continuity of $\gamma_C$
on every $u^{-1}H$ and globally.

The morphism sends $\theta$ to $T$. It is unique because
\eqref{i:normal}, with $T=\theta$, represents every positive-degree
endomorphism. Finally, all Hom spaces in the finite topology are
closed subgroups of products of discrete target modules, defined by
closed $A$-linearity equations, so are separated and complete.
Composition is jointly continuous by controlling finitely many inputs
and their finitely many intermediate images. Statement~\ref{i:zero}
identifies the degree-zero topology. This proves the theorem.
\end{proof}

\begin{theorem}[I3: every regular local surface; Proved]\label{i:noncomplete}
Theorem~\ref{i:surface} holds for every regular noetherian local ring
$A$ of dimension two, with closed-point diagonal algebra $\widehat A$.
The completion required in degree one is specified below.
\end{theorem}
\begin{proof}
Put $B=\widehat A$, $S=A\setminus\{0\}$ and retain
$F,E,E_{\mathfrak p},R_{\mathfrak p},P$ from the preceding proof.
Let $C_{\mathfrak p}$ be the separated completion of
$R_{\mathfrak p}$ in the topology $\tau_{\mathfrak p}$ of
\eqref{i:mixed}, now defined using $A$, and put
$C=\prod_{\mathfrak p}C_{\mathfrak p}$. We will prove
\begin{equation}\label{i:noncompletepieces}
 C_{\mathfrak p}=\Hom_A(E_{\mathfrak p},E),\qquad
 \End_A^1(K)=S^{-1}P\oplus C,\qquad
 \End_A^2(K)=S^{-1}C/S^{-1}B.
\end{equation}

\emph{The completion map has no extra kernel.}
The map $R_{\mathfrak p}\to\Hom_A(E_{\mathfrak p},E)$,
$c\mapsto\delta_{\mathfrak p}c$, is injective: a nonzero $c$ is
a unit times a power of $f_{\mathfrak p}$, and therefore acts
surjectively on $E_{\mathfrak p}$. Its composite with the nonzero
$\delta_{\mathfrak p}$ cannot vanish. Exactness and factoriality
still give the same formula for $\ker\delta_{\mathfrak p}$,
so the induced finite topology on $R_{\mathfrak p}$ is exactly
\eqref{i:mixed}.

For density, take $h:E_{\mathfrak p}\to E$ and finitely many
inputs, contained in the cyclic $A$-module generated by
$e=1/(u f_{\mathfrak p}^n)$, $u\notin\mathfrak p$.
Applying $\Hom_A(A/f_{\mathfrak p}^n,-)$ to the residual resolution
leaves a surjection
\[
 E_{\mathfrak p}[f_{\mathfrak p}^n]
 \xrightarrow{\delta_{\mathfrak p}} E[f_{\mathfrak p}^n],
\]
since $A/f_{\mathfrak p}^n$ has a length-one free resolution.
The source is cyclic over $A_{\mathfrak p}$, generated by $e$.
Choose $c\in A_{\mathfrak p}$ with
$\delta_{\mathfrak p}(ce)=h(e)$. Then
$\delta_{\mathfrak p}c$ agrees with $h$ on all the selected
inputs. The Hom space is separated complete, being closed in a
product of discrete copies of $E$. It is therefore precisely the
separated completion $C_{\mathfrak p}$. Taking products gives
$C=\Hom_A(I,E)$ and dense $P_\tau\subset C$.

The first degree-one summand is unchanged. Dualizing the residual
resolution, now using $D(E)=B$, gives
\[
 0\longrightarrow B\longrightarrow C\longrightarrow D(F)
 \longrightarrow E\longrightarrow0.
\]
The first map sends $b$ to $(b\delta_{\mathfrak p})_{\mathfrak p}$.
Localization proves the last equality of \eqref{i:noncompletepieces};
in particular its entire kernel is $S^{-1}B$. The open-restriction
argument above identifies the open subgroup
$H=\{h\in D(F):h(1)=0\}$ with the topological quotient $C/B$,
and again $D(F)=\bigcup_{u\ne0}u^{-1}H$.

\emph{The universal morphism extends over exactly this completion.}
Let $Q$ be a separated complete target algebra with the prescribed
degree zero $F\times P\times B$ and an element $T$ satisfying the
relations. The map $P\to Q^1$, $c\mapsto e_mTc$, is continuous
for $P_\tau$ by the finite-coordinate and $\mathfrak m^NA$ argument
in the preceding proof. A continuous additive homomorphism is
uniformly continuous, so it extends uniquely to a continuous map
$\beta_Q:C\to Q^1$. This extension is right $P$-linear by
density and continuity. It is also left $B$-linear. Indeed,
$A$-linearity first extends from $P_\tau$ to $C$. For $b\in B$
approximate $b$ by $a_j\in A$ in the $\mathfrak m$-adic topology.
Multiplication by $a_j$ converges pointwise to multiplication by $b$
on $\Hom_A(I,E)$, because every tested value in $E$ is killed
by a power of $\mathfrak m$; multiplication in $Q$ is continuous.
Passing to limits proves the assertion. In particular
\[
 \beta_Q((b\delta_{\mathfrak p})_{\mathfrak p})=b e_mTe_P.
\]
Use the preceding formula for $\alpha_Q:S^{-1}P\to Q^1$, and set
\[
 \gamma_Q([c/u])=\beta_Q(c)Te_\eta u_\eta^{-1}
 \qquad(c\in C,\ u\in S).
\]
Clearing denominators proves independence of representation.
Every element $b/u$ from $S^{-1}B$ maps to
$b e_mTe_PTe_\eta u_\eta^{-1}=0$, by $T^2=0$.
The composition rule is
$\beta(c)\alpha(b/u)=[cb/u]$ for the natural right $P$-action
on $C$; the formulas respect this and all diagonal actions.
Continuity of $\gamma_Q$ on $H=C/B$ follows from the quotient
topology and then extends over each $u^{-1}H$. These maps give the
required continuous algebra homomorphism. Density of $P_\tau$ in
$C$ gives uniqueness: all positive-degree terms are topologically
generated by degree zero and $\theta$. This proves the theorem.
\end{proof}

\paragraph{Remaining gate.}
The presentation is proved in every regular local surface, including
the requested $k[[s,t]]$ case and degree two, under the stated
completion convention. In higher dimension a longer residual
resolution has further intermediate terms. Propositions I4--I7
prepare the proof, and Theorem I8 below identifies all higher relations
in the complete open-left-ideal topology category.
Constructing a residual complex or proving its residue theorem does
not establish that assertion. Yekutieli explicitly specifies
semi-topological rings and a tensor topology
\cite[Definitions 1.2.1 and 1.2.11, pp.~15 and 18]{Yekutieli1992};
his construction is not an opened presentation theorem for the whole
endomorphism algebra. Schenzel's completeness criterion explains why
simply identifying the noncomplete arrow module with
$R_{\mathfrak p}$ would have failed; the dense completion in
Theorem~\ref{i:noncomplete} supplies the missing argument.

\begin{proposition}[I4: adjacent arrows in every dimension; Proved]\label{i:adjacent}
Let $\mathfrak p\subset\mathfrak q$ be an immediate specialization
in a regular noetherian local ring $A$. Put $R=A_{\mathfrak q}$,
$\mathfrak m=\mathfrak qR$, and $T=\widehat{A_{\mathfrak p}}$.
For the residual differential $\delta_{\mathfrak p\mathfrak q}$,
the map
\[
 T\longrightarrow\Hom_A(E_{\mathfrak p},E_{\mathfrak q}),
 \qquad c\longmapsto\delta_{\mathfrak p\mathfrak q}c
\]
is injective and dense. Its induced topology has basis
\[
 \mathfrak m^N R+(\mathfrak pR_{\mathfrak p})^n T,
 \qquad N,n\geq1.
\]
The Hom module is its separated completion for this topology.
\end{proposition}
\begin{proof}
Both hulls are $R$-modules, so $A$-linearity is equivalent to
$R$-linearity. Set $h=\operatorname{ht}(\mathfrak pR)$; then
$\dim R=h+1$. Choose parameters $x_1,\ldots,x_h$ in the regular
local ring $R_{\mathfrak p}$. In its top local cohomology,
which is its injective hull, the cyclic submodules generated by
\[
 e_n=[1/(x_1\cdots x_h)^n],\qquad
 J_n=(x_1^n,\ldots,x_h^n)R_{\mathfrak p},
\]
exhaust $E_{\mathfrak p}$, and $\operatorname{Ann}(e_n)=J_n$.
The ideals $J_n$ are cofinal with powers of $\mathfrak pR_{\mathfrak p}$.
For $h=0$ use $e_n=1$ and $J_n=0$ in the fraction field.
These are the usual Gorenstein and local-cohomology identifications
\cite[Lemmas 47.21.3, 47.21.7, and proof of 47.21.10]{StacksGorenstein}.
Clearing finitely many coefficients over $R_{\mathfrak p}$ shows
that every finite set in $E_{\mathfrak p}$ lies in
$R(e_n/u)$ for suitable $n$ and $u\in R\setminus\mathfrak pR$.

Let $I_n=J_n\cap R$. This ideal is $\mathfrak pR$-primary,
so $R/I_n$ has dimension and depth one. The finite-free-resolution
bound for a regular local ring gives
$\operatorname{pd}_R(R/I_n)\leq h$
\cite[Proposition 10.110.1]{StacksRegularResolution}.
Applying $\Hom_R(R/I_n,-)$ to the residual resolution leaves only
\[
 E_{\mathfrak p}[I_n]\xrightarrow{\delta_{\mathfrak p\mathfrak q}}
 E_{\mathfrak q}[I_n]
\]
in degrees $h,h+1$. The last map is surjective. Its kernel
$W_n=\Ext_R^h(R/I_n,R)$ is finite over $R$.

Given a homomorphism $f$ and inputs in $R(e_n/u)$, lift
$f(e_n/u)$ through this surjection. Since the Artinian complete
intersection $R_{\mathfrak p}/J_n$ is Gorenstein,
$E_{\mathfrak p}[I_n]$ is generated by $e_n/u$ over
$R_{\mathfrak p}$. Thus some $c\in R_{\mathfrak p}$ gives
$\delta_{\mathfrak p\mathfrak q}(ce_n/u)=f(e_n/u)$.
This proves density. Every nonzero $c\in T$ acts surjectively on
$E_{\mathfrak p}$: its Matlis dual is multiplication by $c$ on
the regular complete local domain $T$. Since the displayed
surjection shows $\delta_{\mathfrak p\mathfrak q}\ne0$, the
coefficient map is injective.

Identify $E_{\mathfrak p}[I_n]$ with $(R/I_n)_{\mathfrak p}$
using $e_n$. The finite submodule $W_n$ localizes to this whole
ring, since $E_{\mathfrak q}[I_n]$ vanishes at $\mathfrak p$.
It is therefore commensurable with $R/I_n$: clearing its finite
set of denominators, and then lifting $1$ after localization,
gives $a_n,b_n\in R\setminus\mathfrak pR$ with
\[
 a_n(R/I_n)\subseteq W_n\subseteq b_n^{-1}(R/I_n).
\]
Vanishing on $e_n/u$ means that $c\bmod J_nT$ belongs to
$uW_n$. As $u$ varies, these subgroups are cofinal with
$uR+J_nT$. The ideal $(u,I_n)$ is $\mathfrak m$-primary, and
$\mathfrak m^N$ contains an element outside $\mathfrak pR$.
Thus this family is cofinal with
$\mathfrak m^N R+J_nT$. Cofinality of the $J_n$ with the
$\mathfrak p$-adic powers proves the asserted basis.
Finally the Hom module, being a closed subgroup of a product of
discrete copies of $E_{\mathfrak q}$, is separated complete.
The injective dense map therefore identifies it with the separated
completion just described.
\end{proof}

\begin{proposition}[I5: dense generation in every dimension; Proved]\label{i:dense}
The graded subalgebra of $\End_A^*(K)$ generated by its diagonal
degree-zero algebra and $\theta$ is dense in each degree for the
finite topology.
\end{proposition}
\begin{proof}
First fix $\mathfrak p\subseteq\mathfrak q$. Induct on the height
gap to approximate every map $E_{\mathfrak p}\to E_{\mathfrak q}$
on a finite set by finite sums of saturated paths of residual
arrows with diagonal coefficients. Gap zero follows from
Proposition~\ref{i:zero}; gap one is Proposition~\ref{i:adjacent}.
For larger gap work over $R=A_{\mathfrak q}$, of dimension $d$,
and put the prescribed inputs in a cyclic $R$-submodule $Re$,
where $e=e_n/u$ is as in the preceding proof. Its annihilator $I$
is $\mathfrak pR$-primary. Hence $R/I$ has positive depth, and
\cite[Proposition 10.110.1]{StacksRegularResolution} gives
$\operatorname{pd}_R(R/I)\leq d-1$.
The last differential of $\Hom_R(R/I,K_R)$ is consequently
surjective:
\[
 \bigoplus_{\substack{\operatorname{ht}(\mathfrak rR)=d-1\\
                      \mathfrak p\subseteq\mathfrak r\subset\mathfrak q}}
 E_{\mathfrak r}[I]\longrightarrow E_{\mathfrak q}[I].
\]
Lift $f(e)$; its preimage has finitely many nonzero components
$y_{\mathfrak r}$. The maps $Re\to E_{\mathfrak r}$ sending
$e$ to $y_{\mathfrak r}$ extend to maps
$g_{\mathfrak r}:E_{\mathfrak p}\to E_{\mathfrak r}$ by
injectivity. Thus $\sum_{\mathfrak r}
\delta_{\mathfrak r\mathfrak q}g_{\mathfrak r}$ agrees with
$f$ on $Re$. Approximate each $g_{\mathfrak r}$ on the prescribed
finite inputs by induction, and compose with
$\delta_{\mathfrak r\mathfrak q}$.

For the whole complex, finitely many inputs involve only finitely
many source summands, and their images involve only finitely many
target summands. Apply the preceding construction to these finitely
many component maps, insert endpoint idempotents, and sum. The
result has exactly the prescribed values on the original inputs.
Every path lies in the algebra generated by the diagonal algebra
and $\theta$, proving density.
\end{proof}

\paragraph{Refined higher-dimensional gate.}
Dense generation does not prove a presentation of a completed
algebra. It remains to establish the quotient topology and uniform
continuity needed to extend the universal map from the dense path
subalgebra, and to identify the entire closed ideal of relations
with the ideal generated by $\theta^2$ and incidence.
Yekutieli's Theorem~1.14, p.~9 of the opened author version,
constructs an adelic DG action; Lemma~1.11, pp.~7--8, is an
approximation statement \cite{Yekutieli2003}. An action does not
itself identify its kernel or supply the required quotient topology.
The full presentation in dimension at least three remains
\emph{Undetermined}.

\begin{proposition}[I6: closed-image exactness of residual columns; Proved]\label{i:closedexact}
Let $R$ be a regular noetherian local ring of dimension $d$,
$\mathfrak p$ a nonmaximal prime of height $h$, and $K_R$ its
residual resolution. With finite topologies and differentials by
postcomposition, for $j>h$ one has
\[
 \overline{\operatorname{im}\bigl(
  \Hom(E_{\mathfrak p},K_R^{j-1})\to
  \Hom(E_{\mathfrak p},K_R^j)\bigr)}
 =\ker\bigl(\Hom(E_{\mathfrak p},K_R^j)\to
             \Hom(E_{\mathfrak p},K_R^{j+1})\bigr).
\]
The kernel at $j=h$ is zero.
\end{proposition}
\begin{proof}
Continuity gives the inclusion from left to right. For the other
inclusion, take a cycle $f:E_{\mathfrak p}\to K_R^j$ and put
finitely many prescribed inputs in $Re$ with $e=e_n/u$ as in
Proposition~\ref{i:adjacent}. Its annihilator $I$ is
$\mathfrak p$-primary and
$I R_{\mathfrak p}=(x_1^n,\ldots,x_h^n)$.
The finite module $M=\Ext_R^j(R/I,R)$ vanishes after localization
at $\mathfrak p$ when $j>h$, by the length-$h$ Koszul resolution
there. Thus some $v\notin\mathfrak p$ annihilates $M$.

The element $e/v\in E_{\mathfrak p}$ has the same annihilator
$I$. The restriction of $f$ to $R(e/v)\simeq R/I$ is a cycle.
Its restriction to $Re$ multiplies the cohomology class by $v$,
because $f(e)=v f(e/v)$. This class is zero. Therefore some map
$b:Re\to K_R^{j-1}$ has differential equal to $f$ on $Re$.
Its cyclic source implies that $b$ has finitely many nonzero
components. Extend each component to $E_{\mathfrak p}$ by
injectivity and sum, obtaining
$g:E_{\mathfrak p}\to K_R^{j-1}$. Its differential agrees with
$f$ on the prescribed inputs, proving density in the kernel.

At $j=h$, only the diagonal summand receives a map from
$E_{\mathfrak p}$, so the Hom group is $\widehat{R_{\mathfrak p}}$.
Choose a prime immediately above $\mathfrak p$. The injectivity
in Proposition~\ref{i:adjacent} gives zero kernel for the next
differential.
\end{proof}

For a height-gap-two interval $\mathfrak p\subset\mathfrak q$,
apply this with $R=A_{\mathfrak q}$ and $j=h+1$. It identifies
the kernel of postcomposition from the full adjacent-arrow family
to $\Hom(E_{\mathfrak p},E_{\mathfrak q})$ as the closure of
the family obtained from the preceding residual differential and
$\widehat{A_{\mathfrak p}}$. This is an exact relation statement
in the actual Hom topologies. It does not identify the closure
inside an independently completed free algebra. The quotient-topology
and universal-map continuity gate therefore remains.

\begin{proposition}[I7: strict exactness above the first arrow over a complete ring; Proved]
Let $A$ be a complete regular local ring, let
$E=E_A(A/\mathfrak m)$, and let $\mathfrak p\ne\mathfrak m$.
Write $h=\operatorname{ht}\mathfrak p$ and
$C^j=\Hom_A(E_{\mathfrak p},K_A^j)$ with the finite topology.
The complex $C^\bullet$ is exact as a complex of abstract modules.
Moreover, for every $j>h+1$, its differential induces an open
surjection
\[
 C^{j-1}\longrightarrow Z^j:=\ker(C^j\to C^{j+1}).
\]
No assertion that the first-arrow topology is the local adic
topology is included.
\end{proposition}
\begin{proof}
Matlis' endomorphism theorem identifies
$A=\Hom_A(E,E)$ because $A$ is complete. For a projective
resolution $P_\bullet\to E_{\mathfrak p}$, tensor--Hom adjunction
and exactness of $\Hom_A(-,E)$ give
\[
 \Ext_A^j(E_{\mathfrak p},A)
   =\Hom_A\bigl(\operatorname{Tor}_j^A(E_{\mathfrak p},E),E\bigr).
\]
This uses no finite-generation hypothesis on $E_{\mathfrak p}$:
the adjunction is applied to the whole projective resolution.
Choose $u\in\mathfrak m\setminus\mathfrak p$. The source
$E_{\mathfrak p}$ is an $A_u$-module, whereas $E_u=0$.
Flat localization and derived tensor associativity show
\[
 E_{\mathfrak p}\mathbin{\otimes}_A^{\mathbf L}E
  \simeq
 E_{\mathfrak p}\mathbin{\otimes}_{A_u}^{\mathbf L}E_u=0.
\]
All the displayed Ext groups vanish. Since $K_A^\bullet$ is an
injective resolution of $A$, this proves exactness of $C^\bullet$.

To prove openness, fix a finite set of inputs in $E_{\mathfrak p}$
and put them in a cyclic submodule $N=Ae$ as in
Proposition~\ref{i:adjacent}. Its annihilator $I$ is
$\mathfrak p$-primary and $I A_{\mathfrak p}$ is generated by
powers of a regular system of parameters. If $j>h+1$, the finite
module $\Ext_A^{j-1}(A/I,A)$ vanishes after localization at
$\mathfrak p$. Choose $v\notin\mathfrak p$ annihilating it, and
put $N'=A(e/v)$.

Let $f\in Z^j$ vanish on $N'$. Abstract exactness gives
$g\in C^{j-1}$ with $dg=f$. The restriction of $g$ to $N'$
is a cycle. Restricting further to $N$ multiplies its cohomology
class by $v$, so the class is zero. There is therefore
$b:N\to K_A^{j-2}$ with $db=g|_N$. This map has only finitely
many nonzero injective components; extend each to
$E_{\mathfrak p}$ and sum to obtain $\widetilde b$.
Then $g-d\widetilde b$ vanishes on $N$ and still maps to $f$.
Thus the image of the annihilator of $N$ contains the open
subgroup $Z^j\cap\operatorname{Ann}(N')$. This is exactly
openness of the surjection.
\end{proof}

This removes an abstract exactness and quotient-topology obstruction
for complete-ring columns above their first arrow. It does not
supply the missing uniform construction on arbitrary families of
intermediate primes in an independently defined universal algebra.
In particular it does not promote the full higher-dimensional
presentation to a proved result.


\begin{theorem}[I8: the full presentation in the finite-topology category; Proved]\label{i:linearuniversal}
Let $A$ be a regular noetherian local ring of finite dimension, and
let $K$ be its residual resolution, normalized by height. Put
$B=\prod_{\mathfrak p}\widehat{A_{\mathfrak p}}$.
The finite-topology algebra $\End_A^*(K)$ is universal among
separated complete graded topological $A$-algebras whose topology
has a basis of open left ideals, equipped with the prescribed
continuous degree-zero map from $B$ and an element $T$ of degree one
satisfying
\[
 T^2=0,\qquad e_{\mathfrak q}Te_{\mathfrak p}=0
 \quad\text{unless }\mathfrak p\subset\mathfrak q
 \text{ is an immediate specialization}.
\]
In particular the completed presentation in this category is an
isomorphism in every degree and in every dimension. The universal
property allows a continuous map from $B$; requiring it to identify
degree zero gives the version in the question.
\end{theorem}

\begin{proof}
We first prove a statement about discrete representations that does
not assume the presentation.

\emph{Cousin complexes.}
Call a complex $M$ a Cousin complex here if
\[
 M^i=\bigoplus_{\operatorname{ht}\mathfrak p=i}M_{\mathfrak p},
\]
where $M_{\mathfrak p}$ is an $A_{\mathfrak p}$-module all of whose
elements are killed by powers of $\mathfrak pA_{\mathfrak p}$.
Its differentials are $A$-linear. A map between two such point
modules is zero unless the source prime is contained in the target
prime, by the element-annihilator argument of
Proposition~\ref{i:zero}. Thus the differential has the required
incidence property. The action of
$\widehat{A_{\mathfrak p}}$ on $M_{\mathfrak p}$ is canonical:
on each element it is determined modulo a power of
$\mathfrak pA_{\mathfrak p}$. Every $A$-linear morphism of Cousin
complexes commutes with these completed diagonal actions.

\emph{A universal property of a complete-intersection Cousin
resolution.}
Let $J=(x_1,\ldots,x_h)$ be generated by an $A$-regular sequence,
put $N=A/J$, and put
\[
 C_J=\Hom_A(N,K)=K[J].
\]
The Koszul resolution and its self-duality give
$\Ext_A^i(N,A)=0$ for $i\ne h$ and
$\Ext_A^h(N,A)\simeq N$
\cite[Lemma 48.15.5, tag 0BQZ]{StacksKoszulDuality}.
Choose this last identification, so that
$N$ identifies with $Z^h(C_J)$. Every prime containing $J$ has
height at least $h$, so $C_J^i=0$ for $i<h$. In degree $h$ its
terms identify with
\[
 C_J^h=\bigoplus_{\substack{\mathfrak r\supset J\\
                            \operatorname{ht}\mathfrak r=h}}
 N_{\mathfrak r},
\]
with the augmentation $N\to C_J^h$ given by localization.
Indeed $N$ is Cohen--Macaulay, has no embedded primes, and the
Koszul self-duality identifies the localized top Ext at each
minimal prime with the corresponding Artinian complete intersection.

For every Cousin complex $M$, an $A$-linear map
$N\to Z^h(M)$ extends uniquely to a morphism
\[
                         C_J\longrightarrow M.          \tag{I8.1}
\]
Here is the full extension argument. In degree $h$, localize the
given map at each height-$h$ prime and take their direct sum. Suppose
degrees through $j-1$ have been defined compatibly with the
differentials, where $j>h$. For a prime $\mathfrak q$ of height
$j$, project the composite
$C_J^{j-1}\to M^{j-1}\to M^j$ to $M_{\mathfrak q}$ and
localize at $\mathfrak q$. It kills the kernel of
\[
 (C_J^{j-1})_{\mathfrak q}\longrightarrow
 (C_J^j)_{\mathfrak q}=C_{J,\mathfrak q}^j.
\]
For $j=h+1$ this follows because the prescribed map on $N$ is a
cycle; for larger $j$ it follows from the already established chain
identity and $d_M^2=0$. The displayed arrow is surjective because
the localized complex has no cohomology in degree $j>h$ and no
terms above $j$. Consequently the composite factors uniquely
through $C_{J,\mathfrak q}^j$. Take the direct sum over
$\mathfrak q$. Each source element has finitely many components,
so this gives a well-defined morphism, and comparison of all target
components proves the chain identity. Surjectivity in this
construction also proves uniqueness at every stage. This proves
(I8.1).

\emph{These resolutions generate every Cousin complex.}
Take $v\in M_{\mathfrak p}$ and set
$h=\operatorname{ht}\mathfrak p$. Its differential has finitely
many nonzero components, at primes $\mathfrak q$ strictly containing
$\mathfrak p$. For each such component choose an element of
$\mathfrak q\setminus\mathfrak p$ and a power killing that
component. The product $a$ of these powers is outside
$\mathfrak p$ and satisfies $d_M(av)=0$.
The ideal $\operatorname{Ann}_A(v)$ is $\mathfrak p$-primary:
elements outside $\mathfrak p$ act invertibly, and a power of
$\mathfrak p$ kills $v$. Since $A$ is regular, hence
Cohen--Macaulay, this ideal contains an $A$-regular sequence of
length $h$. Explicitly, after choosing $i<h$ terms, the quotient
is Cohen--Macaulay and its associated primes have height $i$ in
$A$. The height-$h$ annihilator is contained in none of them; prime
avoidance supplies the next nonzerodivisor. This proves the assertion
inductively, using the Cohen--Macaulay quotient property
\cite[Lemma 10.104.2, tag 02JN]{StacksCMRegularSequence}.
Choose such a sequence and its ideal $J$.
The map $N=A/J\to Z^h(M)$ sending $1$ to $av$ extends by
(I8.1) to $C_J\to M$.

The $\mathfrak p$-component $z_{\mathfrak p}$ of the image of
$1\in N$ in $C_J^h$ generates
$C_{J,\mathfrak p}^h\simeq N_{\mathfrak p}$. The extension sends
$z_{\mathfrak p}$ to $av$, and therefore sends
$z_{\mathfrak p}/a$ to $v$. Thus every vector is in the image of
one such morphism. Taking the sum of these maps gives a surjection
of Cousin complexes
\[
                    P:=\bigoplus_\lambda C_{J_\lambda}
                    \longrightarrow M.                  \tag{I8.2}
\]

\emph{All residual endomorphisms act on the quotient.}
Every $C_J=K[J]$ is preserved by all $A$-linear graded
endomorphisms of $K$. Hence $P$ is a discrete continuous
$\End_A^*(K)$-module: a finite set of vectors of $P$ involves only
finitely many inputs in copies of $K$.
The kernel of (I8.2) is preserved by $B$ and by $\theta$ because
the map is a morphism of Cousin complexes. It is consequently
preserved by every word in these operators. It is also preserved by
every $f\in\End_A^*(K)$. To see this without assuming any
relation theorem, take one vector in the kernel, involving
finitely many inputs $z_1,\ldots,z_r$ in copies of $K$.
Proposition~\ref{i:dense} supplies a word polynomial $w$ in $B$
and $\theta$ agreeing with $f$ on all $z_i$. Thus the two operators
give exactly the same value on this vector, and $w$ preserves the
kernel. So $f$ preserves it as well.

It follows that $M$ has a continuous action of the whole
$\End_A^*(K)$ extending its diagonals and differential.
Continuity follows by lifting any finite set of vectors through
(I8.2) and using the corresponding finite set of inputs in $K$.
Uniqueness follows from density. The same density argument shows
that every morphism of Cousin complexes commutes with this action.
This construction proves all relations on every discrete Cousin
representation; it uses density to preserve a kernel in a known
discrete representation, not as a substitute for a presentation.

\emph{Passage to the complete target algebra.}
Let $Q$ be a target in the statement. Replace it by the closed
subalgebra generated by $B$ and $T$. Insert finite sums of the
point idempotents between factors in a word and pass to their
limits. The incidence relations show that a word of degree greater
than $\dim A$ vanishes. Thus this closed subalgebra has finitely
many nonzero graded pieces. Its topology still has a basis of open
left ideals. Replacing an open left ideal $L$ by
$\bigoplus_i(L\cap Q^i)$ gives a smaller open homogeneous left
ideal, since there are finitely many grades. These may therefore
be used as a basis.

For such an ideal $L$, the discrete graded left module $Q/L$
decomposes under $B$ into point-local torsion summands. In fact the
annihilator in $B$ of each vector is open, so it contains the full
factors outside a finite set of primes and suitable powers of the
maximal ideals at the remaining primes. The operator $T$ is its
differential. Decomposing by the constant offset
\emph{graded degree minus prime height} expresses $Q/L$ as a sum
of shifted Cousin complexes. By the result just proved it has a
unique continuous $\End_A^*(K)$-action extending its given
$B,T$-action.

Define
\[
 \phi_L:\End_A^*(K)\longrightarrow Q/L,
 \qquad f\longmapsto f(1+L).
\]
These maps are continuous and compatible as $L$ shrinks, because
the quotient maps are morphisms of the Cousin representations.
Completeness and separatedness give
$Q\simeq\varprojlim_L Q/L$ as topological additive groups.
The $\phi_L$ therefore define a continuous graded map
$\phi:\End_A^*(K)\to Q$, sending $B$ to its prescribed image
and $\theta$ to $T$. It is multiplicative: this holds on word
polynomials, which are dense, and both multiplications and $\phi$
are continuous. Density proves uniqueness as well.

Finally the finite topology on $\End_A^*(K)$ has a basis of open
homogeneous left ideals, namely annihilators of finite sets of
homogeneous inputs. It is separated and complete as proved above.
Thus it belongs to the target category and has exactly the asserted
universal property.
\end{proof}

\paragraph{Completion convention after I8.}
Theorem~\ref{i:linearuniversal} closes every higher-dimensional
relation and quotient-topology issue for the completion whose
neighborhoods are open left ideals, the intrinsic finite topology of
an endomorphism ring. This is stronger than an adelic action or dense
generation statement: it constructs the unique continuous map to
every complete target in that category. It does not identify an
unspecified completion with this one. The transcription of 43:66
says separated complete graded topological algebra but does not
unambiguously state the open-left-ideal convention; 43:64 also
mentions several possibly incompatible linear topologies. If
``topological algebra'' there is instead taken to allow arbitrary
topologies without a basis of open left ideals, that enlarged
universal problem is not settled by I8. The surface proof I2--I3
already works in that enlarged category.


\begin{proposition}[I7b: strict dense column quotients without completeness]
Let $A$ be a regular noetherian local ring and let $K$ be its residual
resolution, normalized by heights. Put
$C_h^j=\operatorname{Hom}_A(K^h,K^j)$ and give it the finite topology.
For $j>h+1$, postcomposition with the differential induces an open
surjection
\[
 C_h^{j-1}\longrightarrow
 \operatorname{im}(d:C_h^{j-1}\to C_h^j),
\]
where the image has its subspace topology. That image is dense in
$Z_h^j=\ker(d:C_h^j\to C_h^{j+1})$.
\end{proposition}
\begin{proof}
First fix a prime $\mathfrak p$ of height $h$ and write
$C_{\mathfrak p}^j=\operatorname{Hom}_A(E_{\mathfrak p},K^j)$.
Finitely many prescribed inputs lie in a cyclic submodule $N=Ae$
whose annihilator $I$ is $\mathfrak p$-primary and whose localization
is generated by powers of a regular system of parameters of
$A_{\mathfrak p}$. Since $j-1>h$, the finite module
$\operatorname{Ext}_A^{j-1}(A/I,A)$ vanishes after localization at
$\mathfrak p$. Choose $v\notin\mathfrak p$ annihilating it and
put $N'=A(e/v)$.

If $f$ belongs to the image of $d$ and vanishes on $N'$, choose an
existing preimage $g$. Then $g|_{N'}$ is a cycle. Restriction to
$N$ multiplies its cohomology class by $v$, so there is
$b:N\to K^{j-2}$ with $db=g|_N$. Its image has finitely many
injective components; extend them to $E_{\mathfrak p}$ to obtain
$\widetilde b$. The map $g-d\widetilde b$ vanishes on $N$ and
still maps to $f$. Therefore
\[
 \operatorname{im}(d)\cap\operatorname{Ann}(N')
       \subseteq d\bigl(\operatorname{Ann}(N)\bigr),
\]
which proves openness. Membership in the image supplied the
preimage, so completeness of $A$ was not used. Density in cycles
is Proposition~I6.

Finally,
\[
 C_h^j=\prod_{\operatorname{ht}\mathfrak p=h}
             C_{\mathfrak p}^j
\]
topologically: a finite set of inputs in $K^h$ meets finitely many
source summands. The product of open surjections is open, and the
product of dense subgroups is dense. This proves the assertion.
\end{proof}

\paragraph{Exact remaining gate.}
If $Q$ is a separated complete topological abelian group and a
continuous additive map $C_h^{j-1}\to Q$ kills $d(C_h^{j-2})$,
then I6, the proposition, and completeness extend it uniquely to
$Z_h^j$. This is a conditional descent result: it requires a
continuous map on the entire preceding-degree family. It neither
constructs that map for arbitrary adjacent-arrow families nor
controls the assembly of endpoint-localized families in a target
without open left ideals. The first-arrow case is excluded, and
the full broad-topology presentation remains undecided.

